\section{A twist preserving the original CM type and its theta line}
\label{sec:source-type-twist}

Put $k=K_d=\mathbb Q(\sqrt{-d})$, where $d>3$ is the
prime in AC-G32, and let $H_d$ be its Hilbert class field.
The letter $H_d$ here denotes a number field, not a source
\texttt{Host} record.  Write $E=E_h$ for the field of AC-G32,
$g=[E:k]=2^{h+2}$, and $\Phi_k$ for the CM type consisting
of all embeddings inducing the chosen embedding of $k$.
Let $AA=\mathfrak b^{-1}$ and $\xi_0=-1/\sqrt{-d}$ be
its original principal CM lattice.  In this section its type
remains $\Phi_k$; the unit modification of AC-G33 is not used.

\begin{theorem}[AC-G35: the original type has a controlled principal twist]
\label{AC-G35}
The complex principally polarized abelian variety defined by
$(E,\Phi_k,AA,\xi_0)$ has a model $(A,\lambda)$ over $H_d$
with all its geometric endomorphisms defined there.
In particular its full $\mathcal O_E$ action is defined there.
The central element $u=[\omega]$ of that action preserves
$\lambda$, although it is not a scalar on the tangent space.
There is a symmetric ample line bundle $\mathcal L$ over $H_d$
inducing $\lambda$, invariant under $u$, with an order-three
linearization that acts as the identity on the fiber at zero.
For this linearization,
\[
             u\bigm|H^0(A,\mathcal L)=1.               \tag{G35.1}
\]

Set $F=EH_d$.  The character $\chi_{269}:G_F\to\mu_3$
defines a twist $(A_\chi,\mathcal L_\chi)$ over $F$,
with the same geometric principal polarization and original
CM type.  With the convention that its descent cocycle acts
on covariant Tate modules by $\chi_{269}$,
\[
 \begin{split}
 T_\ell(A_\chi)&\simeq
      T_\ell(A)\otimes_{\mathcal O_E\otimes\mathbb Z_\ell}
            (\mathcal O_E\otimes\mathbb Z_\ell)(\chi_{269}),\\
 \operatorname{Hom}_{\mathcal O_E\otimes\mathbb Z_\ell}
       (T_\ell(A),T_\ell(A_\chi))
      &\simeq(\mathcal O_E\otimes\mathbb Z_\ell)(\chi_{269})
 \end{split}
                                                        \tag{G35.2}
\]
for every $\ell$, including $\ell=3$.
The dual convention for contravariant $H^1$ reverses the
character.  On the principal theta section line, however,
\[
             H^0(A_\chi,\mathcal L_\chi)
                   \simeq H^0(A,\mathcal L)_F           \tag{G35.3}
\]
with its unchanged descent action.
Thus this particular principal theta line does not acquire
$\chi_{269}$ merely because (G35.2) does.

The model extension has exact cost
\[
 [F:E]=h(k),\quad \mathfrak D_{F/E}=1,\quad
       |D_F|=|D_E|^{h(k)},\quad\operatorname{rd}(F)=\operatorname{rd}(E),
                                                        \tag{G35.4}
\]
where $h(k)$ is odd and $h(k)\le d$.
Every field of definition containing $k$ for this exact
complex torus must contain $H_d$; in particular $E$ itself
cannot be such a field when $h(k)>1$.
These statements concern the model field of the geometric
object.  They do not substitute $F$ for the arithmetic
\texttt{Host.E} in \texttt{NData.sha\_eq}.
\end{theorem}
\begin{proof}
\proofstep{1. An explicit model with all Hom maps.}
The classical input is the elliptic CM theorem:
$k(j(\mathcal O_k))=H_d$, and its Artin class action sends an
ideal lattice $I$ to $\mathfrak a^{-1}I$.
Exact locators are J.~S.~Milne,
\emph{Modular Functions and Modular Forms}, version~1.31,
Theorem~12.18, printed p.~126, and the ideal-isogeny
construction in Lemma~12.19 and Proposition~12.20 on that
same page; \url{https://www.jmilne.org/math/CourseNotes/MF.pdf}.
Choose an elliptic curve $C_0/H_d$ with this $j$ invariant.
Such a Weierstrass model exists over its $j$ field;
here $j\ne0,1728$ because $d>3$.
Its geometric endomorphisms are all defined over $H_d$.
Indeed their derivatives on a differential defined over $H_d$
lie in $k$.  Galois conjugation over $H_d$ fixes those derivatives,
and the derivative map on elliptic endomorphisms is injective
in characteristic zero.  To justify applying Galois here,
a geometric endomorphism is defined over an algebraic extension:
its finite kernel consists of algebraic torsion points, its
quotient is algebraic, and an isomorphism of the resulting
elliptic curves is algebraic.  This proves the claimed descent.

As an $\mathcal O_k$ module, $AA$ is projective of rank $g$.
Write it as a direct sum of invertible ideals
$I_1\oplus\cdots\oplus I_g$.
For each $i$, choose an integral ideal $\mathfrak a_i$ whose
inverse is homothetic to $I_i$.
The quotient
$C_i=C_0/C_0[\mathfrak a_i]$ and its isogeny
$q_i:C_0\to C_i$ are defined over $H_d$, because the ideal
kernel is cut out by endomorphisms already defined there.
Their complex lattices, with one common elliptic period,
are homothetic to the $I_i$.  It follows that
$A=\prod_i C_i$ has the required complex torus.

All geometric Hom maps between these factors are defined
over $H_d$.  For a geometric map $f:C_i\to C_j$, the composite
$q_j^\vee f q_i$ is an endomorphism of $C_0$ defined over $H_d$;
the dual isogeny uses the canonical elliptic polarizations.
The map $f\mapsto q_j^\vee f q_i$ is injective, since both
isogenies are surjective with finite kernels.
Applying every Galois automorphism proves that $f$ is fixed.
Hence every endomorphism of $A$, its specified $\mathcal O_E$
action and its principal polarization map descend to $H_d$.
The complex polarization used here is exactly the trace form
of AC-G32; the audited CM lattice dictionary is Milne's
\emph{Complex Multiplication}, Example~2.9, pp.~25--26,
and Proposition~3.17, p.~31, already used in AM-H17.

\proofstep{2. An actual symmetric principal line bundle.}
This step supplies more than a polarization morphism.
Let $\lambda_0$ be the product of the canonical elliptic
polarizations and write $Q=\lambda_0^{-1}\lambda$ as a
Hermitian matrix of maps between the $C_i$.
Its diagonal entries are integers $q_{ii}$, since the
Rosati-fixed subring of $\mathcal O_k$ is $\mathbb Z$.
Every off-diagonal map is defined over $H_d$ by Step~1.
For $f:C_j\to C_i$, the line bundle
\[
 (\operatorname{pr}_i+f\operatorname{pr}_j)^*
       \mathcal O_{C_i}(0)
 \otimes\operatorname{pr}_i^*\mathcal O_{C_i}(0)^{-1}
 \otimes\operatorname{pr}_j^*f^*\mathcal O_{C_i}(0)^{-1}
                                                        \tag{G35.5}
\]
has zero diagonal polarization entries and cross entries
$f,f^\vee$.  Tensoring these bundles with
$\bigotimes_i\operatorname{pr}_i^*\mathcal O_{C_i}(0)^{q_{ii}}$
therefore constructs a line bundle $\mathcal L_0/H_d$
whose polarization is $\lambda$.
It is symmetric because every factor in (G35.5) is symmetric;
it is ample and principal because its geometric polarization
is the given positive unimodular trace form.
Negative tensor exponents in this construction are harmless:
ampleness is asserted for the final line bundle.

Rosati on $E$ is complex conjugation, so
$u^\dagger u=\overline\omega\omega=1$.
Consequently $u^*\mathcal L_0$ has the same polarization.
Put $D=u^*\mathcal L_0\otimes\mathcal L_0^{-1}$ in
$\operatorname{Pic}^0(A)[2]$; it is two-torsion because both
line bundles are symmetric.
Writing $U=u^*$ additively on this group gives
$(1+U+U^2)D=0$.
On a two-torsion group, $(U-1)UD=D$.
It follows that
\[
                 \mathcal L=\mathcal L_0\otimes U D
                                                        \tag{G35.6}
\]
is $u$-invariant and has the same principal polarization.
All bundles and isomorphisms in this construction are over
$H_d$.  Choose a rigidification at zero.
The isomorphism $u^*\mathcal L\simeq\mathcal L$ can be
uniquely normalized to respect it: two such isomorphisms
differ by a constant in $H_d^\times$.
The product of its three conjugates is an automorphism of
$\mathcal L$ with value one at zero, hence is the identity.
This gives the asserted order-three linearization.

\proofstep{3. The class number and the neutral theta line.}
We give an elementary proof that $h(k)$ is odd.
Conjugation acts on $\operatorname{Cl}(k)$ by inversion.
If a class is fixed, choose a fractional ideal $I$ and
$x\in k^\times$ with $\overline I=(x)I$.
Then $x\overline x$ is a positive rational unit, so is one.
There exists $y\in k^\times$ with $x=y/\overline y$:
take $y=1+x$ when $x\ne-1$, and $y=\sqrt{-d}$ otherwise.
Thus $yI$ is conjugation invariant.
Every conjugation-invariant ideal is a rational fractional
ideal times a product of ramified prime ideals, modulo squares
of those prime ideals.  There is only one ramified prime,
above $d$, and it is the principal ideal $(\sqrt{-d})$.
So the fixed ideal class is trivial.
In particular the finite class group has no two-torsion,
and its order $h(k)$ is odd.

Since $[H_d:\mathbb Q]=2h(k)$ is twice an odd number,
$H_d$ cannot contain both $k$ and the distinct quadratic
field $\mathbb Q(\omega)$.
Thus $\mu_3(H_d)=\{1\}$.
The principal ample line bundle $\mathcal L$ has
$\dim_{H_d}H^0(A,\mathcal L)=1$.
Its normalized $u$ action is an $H_d$-scalar whose cube is one,
so the scalar is one.  This proves (G35.1) for the actual
section line, not merely for an isomorphism class of polarizations.

\proofstep{4. Twist, integral coefficient and local type.}
Over $F=EH_d$, the Kummer character of $\alpha^3=269$
is a homomorphism to the constant group generated by $u$.
Use it as descent data on $(A,\mathcal L)$, with the
linearization of Step~2.  Finite Galois descent is effective
for this projective polarized object, so gives
$(A_\chi,\mathcal L_\chi)/F$.
The cocycle commutes with the $\mathcal O_E$ action and
preserves the principal line bundle.
It leaves the geometric CM type, polarization degree and
stable Faltings height unchanged.

Complex comparison identifies $T_\ell(A)$ with
$AA\otimes\mathbb Z_\ell$ as an
$\mathcal O_E\otimes\mathbb Z_\ell$ module.
It is an invertible module, locally free of rank one,
including at primes above three.
The cocycle acts by multiplication by $\chi_{269}$ on that
module.  Taking the $\mathcal O_E$-linear Hom proves (G35.2)
integrally, without a Fourier idempotent or division by three.
On $H^0(A,\mathcal L)$ the same cocycle acts trivially
by Step~3, proving (G35.3).

At every place where the cocycle is unramified, the twist
becomes isomorphic over an unramified extension.
Thus good ordinary reduction, when present for $A$, is
preserved at such a place.
The character is locally trivial above three by AC-K04,
and is unramified away from the primes above $269$.
This preserves the original type-based ordinary assertions;
it does not claim good reduction for an arbitrary Weierstrass
model at every place.
The new ramification at $269$ is still part of any arithmetic
comparison with the source.

\proofstep{5. Exact model-field cost and necessity.}
The degrees $[E:k]=g$ and $[H_d:k]=h(k)$ are coprime.
Hence $E\cap H_d=k$ and $[EH_d:E]=h(k)$.
The latter extension is unramified, being the base change
of the unramified extension $H_d/k$.
The discriminant tower formula proves (G35.4).

For completeness, the elementary class-number bound needed
here does not require a numerical class calculation.
Minkowski's bound gives an integral ideal of norm at most
$T=(2/\pi)\sqrt d$ in each class.
The number of ideals of any given norm $n$ in a quadratic
field is at most the divisor function $\tau(n)$.
Consequently, since $T>1$,
\[
 h(k)\le\sum_{n\le T}\tau(n)
       \le T(1+\log T)\le T^2<d.                     \tag{G35.7}
\]
The last weak statement $h(k)\le d$ suffices.

To prove the field-of-definition assertion, let $M$ be the
projective $\mathcal O_k$ module underlying $AA$.
Its determinant is an invertible ideal class, its Steinitz class.
By the elliptic CM class action, an Artin class
$[\mathfrak a]\in\operatorname{Gal}(H_d/k)$ sends the
complex product torus to the module $\mathfrak a^{-1}M$.
A complex isomorphism of these tori is an
$\mathcal O_k$-module isomorphism after the common period
is removed; the derivative matrix then has entries in $k$,
because it carries the $k$-span of one lattice to the other.
Taking determinants forces
$[\mathfrak a]^g=1$.
Since $g$ is a power of two and $h(k)$ is odd, this forces
$[\mathfrak a]=1$.
If the torus has a model over a field $L\supset k$,
every element of $G_L$ fixes its complex isomorphism class,
so its image in $\operatorname{Gal}(H_d/k)$ is trivial.
Therefore $H_d\subset L$, as claimed.
\end{proof}

\begin{corollary}[AC-G35-Base: the model field fits the literal auxiliary Base]
\label{AC-G35-Base}
Let $B$ be a source \path{GlobalTheory.Mem.Base} with the
host of AC-G32.  After identifying its coefficient fields
with their chosen complex images, replace
\[
 LL,\ LL_0,\ K_s
       \quad\hbox{by}\quad
 LLH_d,\ LL_0H_d,\ K_sH_d.                            \tag{G35.9}
\]
This defines another literal \texttt{Base} with exactly the
same \texttt{Pre}, ranks and arithmetic \texttt{Host.E}.
The extension of each displayed field is unramified and
has degree at most $h(k)$.  For an old archimedean function
pulled back along restriction of embeddings, the source
normalized average \texttt{meanE} is unchanged.
For old projective integral lattices under literal scalar
extension, their normalized logarithmic index is unchanged
as well.  No assertion of equality with a newly constructed
later theta lattice is included.
\end{corollary}
\begin{proof}
The exact source record is
\path{GlobalTheory.Mem.Base}, in
\path{Definitions/Foundations/RootOfUnityEightQMemCoeffGaloisSubfield.lean},
lines~261--302 in the pinned file retained with the source manifest.
Its additional fields are the number field $LL$ with a chosen
embedding and complex-conjugation involution; a Galois
subfield $LL_0$; a subfield $K_s\subset LL_0$ containing the
host field; specified roots of unity; the clause
\texttt{capit} principalizing each host ideal over $K_s$;
and \texttt{kPsi}, specifying generators of particular
products of conjugate host ideals with absolute-value identities
under every complex embedding.

Both $LL$ and $LL_0$ are Galois over $\mathbb Q$.
The Hilbert class field of an imaginary quadratic field is
also Galois over $\mathbb Q$: conjugation preserves the
unramified abelian extension and acts by inversion on its
Artin class group.  Thus their composita in (G35.9) are
Galois and stable under complex conjugation.
The same complex embedding and its restriction give all
embedding and involution fields of the record.
The field containments are preserved, as are all specified
roots of unity.

An equality of ideals over $K_s$ remains an equality after
extension to $K_sH_d$; extension commutes with products and
integer powers of the nonzero ideals in \texttt{kPsi}.
Its old generator $\beta$ remains nonzero and generates
exactly the extended ideal.
Every complex embedding of $K_sH_d$ restricts to one of
$K_s$, so the old identity for $|\sigma(\beta)|^2$ remains
valid, with the same norm of the original host ideal.
The identical argument preserves \texttt{capit}.
These are all additional fields of the displayed record;
\texttt{Pre} and \texttt{Host.E} were not changed.
All three old fields contain $k$, so their composita with
$H_d$ are unramified extensions of degree at most $h(k)$.

If $m=[LLH_d:LL]$, each old complex embedding has $m$ lifts.
Consequently
\[
 \frac1{[LLH_d:\mathbb Q]}\sum_{\sigma'} f(\sigma'|_{LL})
       =\frac1{[LL:\mathbb Q]}\sum_\sigma f(\sigma),   \tag{G35.10}
\]
which is the exact \texttt{meanE} normalization in
\texttt{DataDIVDefs.lean}.
For an old inclusion of projective integral lattices of the
same rank and finite index, flat base change preserves its
quotient.  Locally a quotient factor $\mathcal O_{LL,v}/\pi^a$
contributes $a\log N v$; above $v$ its contribution is
$a\sum_{w\mid v}e(w/v)\log Nw=ma\log Nv$.
Thus its logarithmic index is multiplied by $m$, and division
by the absolute field degree gives equality.
This does not identify that scalar-extended lattice with
any new lattice defined by subsequent source constructions.
\end{proof}

\paragraph{Quantitative and source interpretation.}
The exact source definitions keep the geometric coefficient
field and the arithmetic host distinct.
In \path{IsDiscriminantTwelfthRootOfIdealLatticeIffDelta1.lean},
\texttt{lamOf} and \texttt{VacIn} use a complex elliptic period
of $K_d$; \texttt{IsDgen} quantifies an auxiliary number field
$M$ containing the embeddings and lattice generators.
In \path{TrueColumnChoicesHeckeCharReferenceEqOne.lean},
\texttt{AmbIn.Ok} asks for that period normalization and
\texttt{IsDgen}, and \texttt{ColIn.A\_AA} retains the original
\texttt{Host.AA0}.
None of these declarations identifies a geometric model field
with \texttt{Host.E}; conversely they do not prove that every
theta lattice and norm estimate is invariant under an arbitrary
change of its coefficient field.

Under the actual window $Y\le x$ and
$\log d\le Y^{\varepsilon/5}(\log Y)^2$, (G35.4)--(G35.7)
give
\[
 \log[F:E]=O(x^{\varepsilon/5}(\log x)^2),\qquad
 [F:\mathbb Q]+\log|D_F|=\exp(o(x^\varepsilon)).        \tag{G35.8}
\]
Here the degree and root discriminant of $E$ are those of
AC-G32, with $2^h$ of the supplied polynomial size in $x$.
Any fixed polynomial in these displayed quantities, multiplied
by a fixed power of $x$, is therefore
$o(\exp(x^\varepsilon/2))$.
This controls precisely such algebraic degree/discriminant costs;
it does not estimate an unlogged theta norm or supply a
Selmer bound over $F$.

Finally, neutral action on the principal line is limited to
(G35.1).  Over $F$, the tangent action of $u$ has eigenvalues
$\omega,\omega^2$, each with multiplicity $g/2$: the original
$k$-induced type restricts equally to the two embeddings of
$\mathbb Q(\omega)$.
Its cotangent symmetric powers and differentiated Fock spaces
can therefore carry nontrivial cubic weights.
Identifying such a charged differentiated space with the
particular source lattice and with the mixed coefficient in
\texttt{LinkN} is a further comparison, not a consequence
of (G35.1) or (G35.2).
The finite descent here uses the same already audited
SGA~1, Expos\'e~V, \S\S4 and~7 dictionary; no new archive
cote is asserted.
